% 扉页与前言
\thispagestyle{empty}
\begin{center}
	\vspace*{3cm}
	{\Huge\heiti __ZH_TITLE__}\\[1.2em]
	{\Large 中文学习版}\\[3em]
	{\large 原著：__AUTHOR__}\\[0.5em]
	{\normalsize \emph{__EN_TITLE__}__EDITION__}\\[\fill]
\end{center}
\clearpage

\chapter*{版权与说明}
\addcontentsline{toc}{chapter}{版权与说明}

\paragraph*{原文} __SOURCE_CITATION__. __LICENSE_SENTENCE__

\paragraph*{本版} 本版是原文的中文学习改编版. 所作的改动:
\begin{enumerate}
	\item 叙述与讲解译为简体中文; 专有名词、定理名与章节标题保留原文的英文写法, 以便与原文对照、做作业;
	\item 在原文较为简略之处增加了补充讲解、例题与推导细节. 凡是新增内容, 一律放在标题为 “\textbf{补充讲解}” 的彩色方框中, 或以 “\textbf{译注}” 脚注的形式出现;
	\item 改正了原文的若干笔误, 每处都以译注注明原文写法;
	\item 版式使用 AJbook 文档类 (李文威, CC BY 4.0), 经由 \url{https://github.com/lichuang/latex-template} 获得.
\end{enumerate}
原作者并未审阅或认可本改编版; 其中的任何错误均属改编者.

\chapter*{阅读说明}
\addcontentsline{toc}{chapter}{阅读说明}
\paragraph*{术语} 书中术语一律使用原文的英文写法 (例如 __TERM_EXAMPLES__), 中文只用于叙述和讲解. 书末索引也按英文术语排列.
\paragraph*{补充讲解} 彩色方框中的内容是中文版新增的, 不属于原文; 考试与作业以原文为准.
